%%%%%%%%%%%%%%%%%%%%%%%%%%%%%%%%%%%%%%%%%%%%%%%%%%%%%%%%%%
%   Autoren:
%   Prof. Dr. Bernhard Drabant
%   Prof. Dr. Dennis Pfisterer
%   Prof. Dr. Julian Reichwald
%%%%%%%%%%%%%%%%%%%%%%%%%%%%%%%%%%%%%%%%%%%%%%%%%%%%%%%%%%

%%%%%%%%%%%%%%%%%%%%%%%%%%%%%%%%%%%%%%%%%%%%%%%%%%%%%%%%%%
%	ANLEITUNG: 
%   1. Ersetzen Sie firmenlogo.jpg im Verzeichnis img
%   2. Passen Sie alle Stellen im Dokument an, die mit 
%      @stud 
%      markiert sind 
%%%%%%%%%%%%%%%%%%%%%%%%%%%%%%%%%%%%%%%%%%%%%%%%%%%%%%%%%%

%%%%%%%%%%%%%%%%%%%%%%%%%%%%%%%%%%%%%%%%%%%%%%%%%%%%%%%%%%
%	ACHTUNG: 
%   Für das Erstellen des Literaturverzeichnisses wird das 
%   modernere Paket biblatex in Kombination mit biber 
%   verwendet - nicht mehr das ältere Paket BibTex!
%
%   Bitte stellen Sie Ihre TeX-Umgebung entsprechend ein (z.B. TeXStudio): 
%   Einstellungen --> Erzeugen --> Standard Bibliographieprogramm: biber
%%%%%%%%%%%%%%%%%%%%%%%%%%%%%%%%%%%%%%%%%%%%%%%%%%%%%%%%%%

\documentclass[fontsize=11pt,BCOR=5mm,DIV=13,parskip=half,listof=totoc,
               paper=a4,toc=bibliography,pointlessnumbers]{scrreprt}



               
\usepackage[utf8]{inputenc}

%% LANGUAGE SETTINGS
%
% @stud: Sprache ggf. anpassen
%
\usepackage[english]{babel}
\usepackage{csquotes} 	% correct quoting using \enquote{}

%%%%%%%%%%%%%
%% ZITIERSTIL
%%%%%%%%%%%%%
%
% @stud: Zitierstil in package biblatex unten wählen
%
% NUMERIC Style - e. g. [12]
% style=numeric 
%
% IEEE Style - numeric kind of style 
% style=ieee 
%
% ALPHABETIC Style - e. g. [AB12]
% style=alphabetic 
%
% HARVARD Style 
% style=apa 
%
% CHICAGO Style 
%style=authoryear
%
% Position des Zitats:
%
% autocite=inline 
%
% (!!) FOOTNOTE POSITION NOT RECOMMENDED IN MINT DOMAIN:
% autocite=footnote
%
\usepackage[backend=biber, autocite=inline, style=authoryear]{biblatex} 	
\usepackage{makeidx}                  % allows index generation
\usepackage{listings}	                %Format Listings properly
\usepackage{lipsum}                   % Blindtext
\usepackage{graphicx}                 % use various graphics formats
\usepackage[english]{varioref}         % nicer references \vref
\usepackage{caption}	                % better Captions
\usepackage{booktabs}                 % nicer Tabs
\usepackage[hidelinks=true]{hyperref} % keine roten Markierungen bei Links
\usepackage{fnpct}                    % Correct superscripts 
\usepackage{calc}                     % Used for extra space below footsepline, in particular
\usepackage{array}
\usepackage{acronym}
\usepackage{algorithm}
\usepackage{algpseudocode}
\usepackage{setspace}
\usepackage{tocloft}
\usepackage[T1]{fontenc}

% Definitionen und Commands
\newcommand{\indextype}{numeric}
\newcommand{\abs}{\par\vskip 0.2cm\goodbreak\noindent}
\newcommand{\nl}{\par\noindent}
\newcommand{\mcl}[1]{\mathcal{#1}}
\newcommand{\nowrite}[1]{}
\newcommand{\NN}{{\mathbb N}}
\newcommand{\imagedir}{img}
\newcommand{\TitelDerArbeit}[1]{\def\DerTitelDerArbeit{#1}\hypersetup{pdftitle={#1}}}
\newcommand{\AutorDerArbeit}[1]{\def\DerAutorDerArbeit{#1}\hypersetup{pdfauthor={#1}}}
\newcommand{\Firma}[1]{\def\DerNameDerFirma{#1}}
\newcommand{\Kurs}[1]{\def\DieKursbezeichnung{#1}}
\newcommand{\Abteilung}[1]{\def\DerNameDerAbteilung{#1}}
\newcommand{\Studiengangsleiter}[1]{\def\DerStudiengangsleiter{#1}}
\newcommand{\WissBetreuer}[1]{\def\DerWissBetreuer{#1}}
\newcommand{\FirmenBetreuer}[1]{\def\DerFirmenBetreuer{#1}}
\newcommand{\Bearbeitungszeitraum}[1]{\def\DerBearbeitungszeitraum{#1}}
\newcommand{\Abgabedatum}[1]{\def\DasAbgabedatum{#1}}
\newcommand{\Matrikelnummer}[1]{\def\DieMatrikelnummer{#1}}
\newcommand{\Studienrichtung}[1]{\def\DieStudienrichtung{#1}}
\newcommand{\ArtDerArbeit}[1]{\def\DieArtDerArbeit{#1}}
\newcommand{\Literaturverzeichnis}{Literaturverzeichnis}

% Page Layout
%\oddsidemargin=0mm
%\evensidemargin=0mm
%\textwidth=159mm
%\topmargin=-18mm
%\headsep=10mm
%\textheight=251mm
%\footheight=15mm

\makeindex

%%%%%%%%%%%%%%%%%%%%%%%%%%%%%%%%%%%
% LITERATURVERZEICHNIS
% @stud: Literaturverzeichnis in Datei bibliography.bib anpassen. 
%
% [Alternative zu Verwendung von \initializeBibliography: Citavi ... (dazu eigenes LaTex Coding verwenden)]
%
\addbibresource{bibliography.bib}
\DefineBibliographyStrings{ngerman}{andothers = {{et\,al\adddot}},}

% Elementare Konfigurationen und Definitionen werden geladen 
% @stud: gegebenenfalls anpassen
%
% !TEX root =  master.tex

%%%%%%%%%%%%%%%%%%%%%%%%%%%%%%%%%%%%%%%%%%%%%%%%%%%%%%%%%%%%%%%%%%
%	ANLEITUNG: 
% Passen Sie gegebenenfalls alle Stellen im Dokument an, die mit 
% @stud 
% markiert sind.
%%%%%%%%%%%%%%%%%%%%%%%%%%%%%%%%%%%%%%%%%%%%%%%%%%%%%%%%%%%%%%%%%%

%%
%% @stud
%%
%% LANGUAGE SETTINGS
\usepackage[english]{babel} 	        % german language
\usepackage{csquotes} 
%\usepackage[english]{babel}          % english language
%\usepackage{csquotes} 	              % correct quoting using \enquote{}

\usepackage{makeidx}                  % allows index generation
\usepackage{listings}	                %Format Listings properly
\usepackage{lipsum}                   % Blindtext
\usepackage{graphicx}                 % use various graphics formats
\usepackage[english]{varioref}         % nicer references \vref
\usepackage{caption}	                % better Captions
\usepackage{booktabs}                 % nicer Tabs
\usepackage[hidelinks=true]{hyperref} % keine roten Markierungen bei Links
\usepackage{fnpct}                    % Correct superscripts 
\usepackage{calc}                     % Used for extra space below footsepline, in particular
\usepackage{array}
\usepackage{acronym}
\usepackage{algorithm}
\usepackage{algpseudocode}
\usepackage{setspace}
\usepackage{tocloft}
\usepackage{rotating}

%% Schriftarten- und Zeichenpakete
\usepackage[T1]{fontenc}
\usepackage[utf8]{inputenc}


%%
%% @stud
%%
%%	FONT SELECTION: Schriftarten und Schriftfamilie
%%%%%%%%%%%%%
%% SCHRIFTART
%%%%%%%%%%%%%
% 0) without decomment: normal font families 
% ...
% 1) Latin Modern 
\usepackage{lmodern}        
% 2) Times 
%\usepackage{mathptmx}         
% 3) Helvetica
%\usepackage[scaled=.92]{helvet} 
%%%%%%%%%%%%%%%%%%
%%	SCHRIFTFAMILIE
%%%%%%%%%%%%%%%%%%
% ohne Serifen
\renewcommand*{\familydefault}{\sfdefault}
\addtokomafont{disposition}{\sffamily}
%
% mit Serifen
%\renewcommand*{\familydefault}{\rmdefault}
%\addtokomafont{disposition}{\rmfamily}
%
% Typewriter
%\renewcommand*{\familydefault}{\ttdefault}
%\addtokomafont{disposition}{\ttfamily}

%%
%% @stud
%%
%% Uncomment the following lines to support hard URL breaks in bibliography 
%\apptocmd{\UrlBreaks}{\do\f\do\m}{}{}
%\setcounter{biburllcpenalty}{9000}% Kleinbuchstaben
%\setcounter{biburlucpenalty}{9000}% Großbuchstaben

%%
%% @stud
%%
%% FOOTNOTES: Count footnotes over chapters
\counterwithout{footnote}{chapter}

%	ACRONYMS
\makeatletter
\@ifpackagelater{acronym}{2015/03/20}
{\renewcommand*{\aclabelfont}[1]{\textbf{{\acsfont{#1}}}}}{}
\makeatother

%	LISTINGS
% @stud: ggf. Namen/Text anpassen (englisch)
\renewcommand{\lstlistingname}{Source Code} 
\renewcommand{\lstlistlistingname}{List of Source Codes}
\lstset{numbers=left,
	numberstyle=\tiny,
    breaklines=true,
	basicstyle=\ttfamily\small}

%	ALGORITHMS
% @stud: ggf. Namen/Text anpassen (englisch)
\renewcommand{\listalgorithmname}{Algorithmenverzeichnis}
\floatname{algorithm}{Algorithmus}

%	PAGE HEADER / FOOTER
%	Warning: There are some redefinitions throughout the master.tex-file!  DON'T CHANGE THESE REDEFINITIONS!
\RequirePackage[automark]{scrlayer-scrpage}
%alternatively with separation lines: \RequirePackage[automark,headsepline,footsepline]{scrlayer-scrpage}

\renewcommand{\chaptermarkformat}{}
\RedeclareSectionCommand[beforeskip=0pt]{chapter}
\clearpairofpagestyles

%\ifoot[\rule{0pt}{\ht\strutbox+\dp\strutbox}DHBW Mannheim]{\rule{0pt}{\ht\strutbox+\dp\strutbox}DHBW Mannheim}
\ofoot[\rule{0pt}{\ht\strutbox+\dp\strutbox}\pagemark]{\rule{0pt}{\ht\strutbox+\dp\strutbox}\pagemark}
\ohead{\headmark}

\providecommand{\TitelDerArbeit}[1]{\def\DerTitelDerArbeit{#1}\hypersetup{pdftitle={#1}}}
\providecommand{\AutorDerArbeit}[1]{\def\DerAutorDerArbeit{#1}\hypersetup{pdfauthor={#1}}}
\providecommand{\Firma}[1]{\def\DerNameDerFirma{#1}}
\providecommand{\Kurs}[1]{\def\DieKursbezeichnung{#1}}
\providecommand{\Abteilung}[1]{\def\DerNameDerAbteilung{#1}}
\providecommand{\Studiengangsleiter}[1]{\def\DerStudiengangsleiter{#1}}
\providecommand{\WissBetreuer}[1]{\def\DerWissBetreuer{#1}}
\providecommand{\FirmenBetreuer}[1]{\def\DerFirmenBetreuer{#1}}
\providecommand{\Bearbeitungszeitraum}[1]{\def\DerBearbeitungszeitraum{#1}}
\providecommand{\Abgabedatum}[1]{\def\DasAbgabedatum{#1}}
\providecommand{\Matrikelnummer}[1]{\def\DieMatrikelnummer{#1}}
\providecommand{\Studienrichtung}[1]{\def\DieStudienrichtung{#1}}
\providecommand{\ArtDerArbeit}[1]{\def\DieArtDerArbeit{#1}}
\providecommand{\Literaturverzeichnis}{Literaturverzeichnis}

\newcommand{\settingBibFootnoteCite}{
	\setlength{\bibparsep}{\parskip}		  % Add some space between biblatex entries in the bibliography
	\addbibresource{bibliography.bib}	    % Add file bibliography.bib as biblatex resource
	\DefineBibliographyStrings{english}{andothers = {{et\,al\adddot}},}
}

\newcommand{\setTitlepage}{
	% !TEX root =  master.tex
% @stud: ggf. Namen/Text anpassen (englisch)
\begin{titlepage}
    \begin{minipage}{\textwidth}
            \vspace{-2cm}
            \noindent \includegraphics[width=0.3\textwidth]{\imagedir/firmenlogo.jpg} \hfill 
            \includegraphics[width=0.2\textwidth]{\imagedir/logo.jpg}
    \end{minipage}
    \vspace{1em}
    %\sffamily
    \begin{center}
        {\textsf{\textbf{\large{\DieArtDerArbeit}}}}\\[6mm]
        {\textsf{\textbf{\Large{}\DerTitelDerArbeit}}} \\[1.5cm]
        {\textsf{\textbf{\large{}Course of Study: Business Informatics}}}\\[6mm]
        {\textsf{\textbf{Specialization: \DieStudienrichtung}}}\vspace{10em}
        
        \begin{minipage}{\textwidth}
            \begin{tabbing}
            Scientific Supervisor: \hspace{0.85cm}\=\kill
            Author: \> \DerAutorDerArbeit \\[1.5mm]
            Company: \> \DerNameDerFirma  \\[1.5mm]
            Department: \> \DerNameDerAbteilung \\[1.5mm]
            Supervisor: \> \DerFirmenBetreuer \\[1.5mm]
            Company: \> SAP Labs India \\[1.5mm]
            Department: \> MD's office \\[1.5mm]
            \end{tabbing}
        \end{minipage}
        
    \end{center}
    \end{titlepage}
	\pagenumbering{roman} % Römische Seitennummerierung
	\normalfont	
}

\newcommand{\initializeText}{
	\clearpage
	\ihead{\chaptername~\thechapter} % Neue Header-Definition
	\pagenumbering{arabic}           % Arabische Seitenzahlen
}

\newcommand{\initializeBibliography}{
	\ihead{}
	\printbibliography[title=\Literaturverzeichnis] 
	\cleardoublepage
}

\newcommand{\initializeAppendix}{
	\appendix
  \ihead{}
  \cftaddtitleline{toc}{chapter}{Appendix}{}
}

% @stud
%
% PERSÖNLICHE ANGABEN (BITTE VOLLSTÄNDIG EINGEBEN zwischen den Klammern: {...})
%
\ArtDerArbeit{Whitepaper} 
\TitelDerArbeit{Exploring Voice-Controlled Drone Technology: PoC Development for SAP Labs India}
\AutorDerArbeit{Sean Tyler Straub}
\Abteilung{Corporate IT}
\Firma{Freudenberg \& Co. KG}
\Studienrichtung{Software Engineering}
\FirmenBetreuer{Kishor Ingale}

\begin{document}

\setTitlepage


%	Kurzfassung / Abstract
% @stud: abstract.tex bearbeiten
\section*{Abstract}
This paper presents the development of two \acp{PoC} for the \ac{ASTRID} project, 
an initiative at SAP Labs India aimed at leveraging \acp{LLM} for voice-controlled drone operations. 

The first \ac{PoC} explores the application of few-shot prompting to generate precise drone command sequences, demonstrated through a PyGame simulation environment. 
It highlights the iterative refinement of prompts to handle varying complexity in user inputs, enabling intuitive and accurate drone behavior.

The second \ac{PoC} addresses real-time video streaming challenges, transitioning from high-latency YouTube-based streaming to an NGINX \ac{RTMP} server. 
This approach achieved a significant reduction in latency to approximately 1-2 seconds, enhancing the responsiveness and usability of drone video feeds for real-world applications.

Together, these \acp{PoC} showcase the potential of integrating \acp{LLM} with drone technology to enable seamless voice control and improved real-time feedback. 
The findings lay the groundwork for further advancements in voice-controlled drone systems, particularly in applications like surveillance, remote monitoring, and emergency response. 
\cleardoublepage


%%%%%%%%%%%%%%%%%%%%%%%%%%%%%%%%%%%

%%%%%%%%%%%%%%%%%%%%%%%%%%%%%%%%%%%
% VERZEICHNISSE und ABSTRACT
%
% @stud: ggf. nicht benötigte Verzeichnisse auskommentieren/löschen
%
\tableofcontents
\cleardoublepage

% Abbildungsverzeichnis
%\phantomsection
%\addcontentsline{toc}{chapter}{\listfigurename}
%\listoffigures
%\cleardoublepage

%	Tabellenverzeichnis
%\phantomsection
%\addcontentsline{toc}{chapter}{\listtablename}
%\listoftables
%\cleardoublepage

%	Listingsverzeichnis / Quelltextverzeichnis
%\lstlistoflistings
%\cleardoublepage

% Algorithmenverzeichnis
% \listofalgorithms
% \cleardoublepage

% Abkürzungsverzeichnis
% @stud: acronyms.tex bearbeiten
% !TEX root =  master.tex
\clearpage
\chapter*{Acronyms}	
\addcontentsline{toc}{chapter}{Acronyms}

\begin{acronym}[XXXXXXX]	
    \acro{ASTRID}{Autonomous System for Remote and Intelligent Drone operations}
    \acro{HLS}{HTTP Live Streaming}
    \acro{LLM}{Large Language Model}
    \acro{PoC}{Proof of Concept}
    \acro{RTMP}{Real-Time Messaging Protocol}

\end{acronym} 
\cleardoublepage

\onehalfspacing

\initializeText

%%%%%%%%%%%%%%%%%%%%%%%%%%%%%%%%%%%%%%%%%%%%%%%%%%%%%%%%%%%%%%%%%%%%%%%%%%%%%%%%%%%%%%%%%%
% KAPITEL UND ANHÄNGE
%
% @stud:
%   - nicht benötigte: auskommentieren/löschen
%   - neue: bei Bedarf hinzufügen mittels input-Kommando an entsprechender Stelle einfügen
%%%%%%%%%%%%%%%%%%%%%%%%%%%%%%%%%%%%%%%%%%%%%%%%%%%%%%%%%%%%%%%%%%%%%%%%%%%%%%%%%%%%%%%%%%

%%%%%%%%%%%%%%%%%%%%%%%%%%%%%%%%%%%
% KAPITEL
%
% @stud: einzelne Kapitel bearbeiten und eigene Kapitel hier einfügen
%

\chapter{Introduction}
\label{introduction}
\nocite{*}

\section{Background}
This document details the development of two \acp{PoC} as part of the \ac{ASTRID} project at SAP Labs India. \ac{ASTRID}, an initiative driven by the Innovations Team, 
seeks to explore the potential of \acp{LLM} in enabling seamless and intuitive voice-based control for drones. 

Modern drones offer sophisticated capabilities, yet their control interfaces often require significant expertise, limiting their accessibility for non-technical users. 
By leveraging advancements in \acp{LLM}, the \ac{ASTRID} project aims to simplify drone operation by enabling natural language voice commands. This vision has applications across domains such as surveillance, remote monitoring, and emergency response.

The project's practical implementation was tested on a DJI Mini 3 Pro drone, a compact and versatile drone known for its advanced capabilities, including high-quality video capture and stability. 
The drone connects to a dedicated controller, which is paired with a mobile device running the DJI app. This setup provides both control inputs and real-time video feedback, forming the base infrastructure for the integration of voice-based operations and low-latency streaming solutions.

The \acp{PoC} developed within this project aim to address key challenges in achieving this vision, laying the groundwork for future enhancements in voice-controlled drone technology.

\section{Preread}
To ensure a robust understanding of the technical aspects required for this project, I reviewed foundational materials, including the DJI user manual and the DJI SDK documentation. 
These resources provided critical insights into drone hardware capabilities, control systems, and the SDK features available for integration. 

In addition to official documentation, I consulted a range of blog posts and technical articles exploring \ac{RTMP} streaming optimization. 
These secondary resources helped shape the approach to minimizing latency and enhancing the real-time responsiveness of drone video feeds. 
This preparatory work informed the design and implementation of the \acp{PoC}, ensuring alignment with the project's goals.

\cleardoublepage

\chapter{Implementation}
\label{implementation}

The \ac{ASTRID} project focuses on enabling drone control through voice commands powered by Large Language Models (\acp{LLM}). 
As part of the initiative, two Proofs of Concept (\acp{PoC}) were developed, each addressing a unique challenge in the application of \acp{LLM} for drone operations.

\section{Methodology: Drone Command Generation}
\subsection{Overview}
The first \ac{PoC} aimed to explore the potential of few-shot prompting to generate a sequence of precise drone commands based on natural language instructions. 
This methodology involved designing a set of adaptable prompts capable of handling varying levels of complexity in user inputs.

\subsection{Prompt Engineering}
Prompt engineering focused on defining clear parameters for movement and rotation. The drone commands followed predefined formats such as "Go Front", "Turn Left", "Go Right", and "Go Up", 
with the goal of translating user requests for movement into a chain of executable instructions. 
Key challenges included refining the prompts to ensure that the \ac{LLM} generated commands that were both accurate and feasible, given the constraints of the drone SDK.

For instance, a simple prompt like "Move forward 50 cm and then turn left 90 degrees" would result in a sequence of commands like:  
\texttt{Go Front, Go Front, Go Front, Go Front, Go Front, Turn Left, Turn Left, Turn Left}.  
This approach allowed the \ac{LLM} to handle both simple and more complex instructions, including geometric shapes and multi-step movements. 
For complex shapes like squares or circles, the prompt was adjusted to guide the \ac{LLM} step-by-step, ensuring that it could break down larger tasks into smaller, actionable commands.

\section{Methodology: Real-time Video Streaming}
\subsection{Overview}
The second \ac{PoC} focused on improving the latency of real-time video streaming from the drone. 
Prior to this work, the Innovations Team had relied on YouTube to stream drone footage, but this solution suffered from significant latency, often exceeding 10 seconds. 
To address this challenge, an NGINX \ac{RTMP} server was set up to directly receive the \ac{RTMP} stream from the drone, offering a significant improvement in latency.

\subsection{RTMP Server Configuration}
The implementation of the \ac{RTMP} server involved several key steps:
\begin{itemize}
    \item Configuring an NGINX server with the \ac{RTMP} module to accept and forward \ac{RTMP} streams.
    \item Setting up the server on a local machine (MacBook) and testing its performance with tools like \texttt{ffmpeg} to measure latency.
    \item Fine-tuning parameters to minimize delay and ensure smooth streaming.
\end{itemize}

\section{Implementation: Drone Command Generation}
\subsection{Few-shot Prompting in PyGame}
The first \ac{PoC} was implemented in a simulation environment using PyGame, where the \ac{LLM}-generated commands could be executed in real-time. 
The simulation employed the same command set as the drone SDK, including basic movement and rotation instructions. 
The prompts were designed to ensure that the \ac{LLM} could generate commands that were both feasible and precise, considering the drone's movement constraints.

The process involved sending user input to the \ac{LLM} via a chatbot interface, where the \ac{LLM} would respond with a series of commands. 
These commands were then executed in the simulation, with the drone moving or rotating based on the generated instructions. 
Prompt engineering was iteratively refined based on feedback from initial tests, ensuring that the \ac{LLM} understood the task and produced relevant commands.

This implementation enabled testing of complex scenarios, such as multiple turns or navigating specific shapes, and helped refine the \ac{LLM}'s response to these instructions.

\subsection{Prompt Engineering and Refinement}
Initially, the \ac{LLM} used for few-shot prompting was GPT-4o-mini, but it was quickly replaced by GPT-4o due to the former's inability to generate accurate drone commands. 
Additionally, system prompts were refined to ensure that the \ac{LLM} could produce precise and actionable instructions for drone movement and behavior. 
The final iteration of the system prompt, as shown in the code snippet \ref{lst:prompt_engineering}, ensured the required accuracy.

\section{Implementation: Real-Time Video Streaming}
\subsection{NGINX RTMP Server for Video Streaming}
The second \ac{PoC} aimed to improve the real-time video streaming experience by reducing latency, transitioning from YouTube to a local NGINX \ac{RTMP} server. 
This solution was tested with the different encoding settings H.264 and H.265 but the results were inconclusive in determining which encoding performed better. 
However, a key configuration, \texttt{sync}, was crucial for ensuring stable streaming. After testing, it was found that a \texttt{sync} value of 5ms provided the best results, 
though further testing under varying conditions is recommended.

The setup process involved installing and configuring the NGINX \ac{RTMP} module on the MacBook, adjusting settings to handle the stream efficiently. 
Results showed that the \ac{RTMP} server reduced latency significantly, achieving a stream delay of approximately 1-2 seconds—substantially better than the previous YouTube-based solution.

Although this setup worked well in the test environment using a mobile hotspot, subsequent tests on SAP's Wi-Fi network were hindered by firewall issues. 
This could be resolved in a production environment by implementing a dedicated router or an alternative network configuration.

\section{Challenges and Limitations}
Several challenges were encountered during the implementation of both \acp{PoC}. For the few-shot prompting task, initial testing with GPT-4o-mini proved insufficient for generating accurate drone commands. 
Transitioning to GPT-4o was essential for achieving the desired results, as it offered improved comprehension and response accuracy. 
Additionally, the system prompts required several iterations to ensure the \ac{LLM} generated precise, actionable instructions for drone movement.

For the \ac{RTMP} streaming, the primary hurdle was the network setup. Although the mobile hotspot provided a temporary solution, a more stable network infrastructure is necessary for long-term deployment. 
Additionally, the tests to determine the best video settings were inconclusive, warranting further investigation to identify the optimal configuration.

Despite these challenges, 
both \acp{PoC} demonstrated the potential of using \acp{LLM} in drone control and highlighted the importance of real-time video streaming in applications such as surveillance and remote monitoring.

\cleardoublepage

\chapter{Conclusion}
\label{conclusion}


\section{Findings}
The \ac{ASTRID} project successfully explored the potential of integrating Large Language Models (\acp{LLM}) for drone control, specifically through voice commands. 
The two Proofs of Concept (\acp{PoC}) demonstrated promising results in both drone command generation and real-time video streaming.

The first \ac{PoC} showcased the feasibility of using few-shot prompting to generate precise drone commands based on natural language instructions. 
The prompt engineering process, combined with iterative refinements, enabled the \ac{LLM} to produce feasible and actionable commands, even for complex movements. 
This work opens up possibilities for more intuitive human-drone interaction through voice control, allowing users to issue commands in a natural and flexible manner.

The second \ac{PoC} addressed the challenge of real-time video streaming from the drone. 
By transitioning from YouTube to a locally hosted NGINX \ac{RTMP} server, significant improvements in latency were achieved, reducing the delay to just 1-2 seconds. 
This is a crucial advancement for applications where real-time video feedback is essential, such as in surveillance and remote monitoring.

\section{Next Steps}
Despite the successes, several areas remain for further development and refinement. The next steps for the \ac{ASTRID} project include:

\begin{itemize}
    \item \textbf{Enhancing Command Generation}:
    While the current \ac{LLM}-generated command sequences are functional, further refinement is needed to handle even more complex scenarios, such as unpredictable environmental conditions or emergency maneuvers. 
    This can be achieved by introducing additional context into the prompts and expanding the command set.

    \item \textbf{Utilizing and Deploying Command Generation in Real-world Scenarios}:
        Testing and deploying the command generation system in real-world scenarios will be crucial to validate its performance and reliability. 
        This involves conducting field tests where the system is used in various operational environments to ensure it can handle diverse and unpredictable conditions. 
        This includes testing under various weather conditions, with different drone models, and in complex operational settings like surveillance and delivery. 
        Feedback from these tests will be invaluable for refining the system and addressing any limitations or challenges that arise.
    
    \item \textbf{Integration with other Systems}:
    Future work should include integrating the \ac{ASTRID} system with other drone control frameworks and software platforms, enabling seamless interoperability across different devices and use cases. 
    This may involve supporting a wider range of drone models, extending the system's compatibility beyond just DJI drones.
    
    \item \textbf{Advanced Real-time Video Streaming}:
    Although the NGINX \ac{RTMP} server reduced latency majorly, further testing is needed to determine the optimal encoding settings and server configurations for varied network conditions. 
    Additionally, evaluating the impact of different hardware configurations, such as dedicated servers or edge devices, will help improve streaming performance in larger-scale applications.
    
    \item \textbf{Exploring Edge Computing for Low-Latency Processing}:
    To further reduce latency, exploring edge computing solutions for processing drone commands and video streams could be beneficial. 
    By offloading some processing to local devices, it may be possible to achieve even lower latency and improve the overall responsiveness of the system.

    \item \textbf{Transcoding from RTMP to HLS}:
    Another potential enhancement is the ability to trans-code \ac{RTMP} streams to \ac{HLS} using the NGINX server. 
    This would enable more flexible and adaptive streaming options, as \ac{HLS} is widely supported across various devices and platforms. 
    Implementing this feature would involve configuring the NGINX server to transcode the incoming \ac{RTMP} stream into HLS segments, 
    allowing for adaptive bitrate streaming and improved compatibility with different viewing devices.

\end{itemize}

In conclusion, the \ac{ASTRID} project lays a strong foundation for voice-controlled drone operations and real-time video streaming, opening the door to more intuitive and efficient drone applications. 
As the system continues to evolve, its potential for use in industries like surveillance, remote inspection, and emergency response holds great promise. 
With further development and real-world testing, \ac{ASTRID} could lead to a new generation of autonomous drones that respond seamlessly to voice commands, with minimal latency and enhanced operational capabilities.

\cleardoublepage
%%%%%%%%%%%%%%%%%%%%%%%%%%%%%%%%%%%

%%%%%%%%%%%%%%%%%%%%%%%%%%%%%%%%%%%
% ANHÄNGE
%
% @stud: einzelne Anhänge bearbeiten und eigene Anhänge hier einfügen 
%        die nachfolgenden Zeilen deaktivieren, wenn keine Anhänge verwendet werden
%
\initializeAppendix
\chapter{Python Code Snippet}
\label{appendix}


\begin{lstlisting}[caption={Prompt Engineering Drone Command Generation}, label={lst:prompt_engineering}]
response = client.chat.completions.create(
    model="gpt-4o",
    temperature=0,
    messages=[
        {
            "role": "system",
            "content": (
                "You are a drone controller in a simulation. You are only allowed to respond using specific commands in the following format: "
                "Go Front, Go Back, Go Left, Go Right, Go Up, Go Down, Turn Right, Turn Left, Land, Return to Home."
                "\n\n"
                "General Instructions:\n"
                "1. Movement commands (Go Front, Go Back, Go Left, Go Right, Go Up, Go Down) move exactly 10 cm per command.\n"
                "2. Rotation commands (Turn Right, Turn Left) rotate the drone by exactly 30 degrees per command.\n"
                "3. 'Land' is used to bring the drone back to the ground (z-coordinate = 0).\n"
                "4. 'Return to Home' commands the drone to return to its starting point.\n"
                "5. Complex figures such as squares, spirals, and zigzags should be broken down into small, clear steps of movement and rotation, ensuring precise replication of the requested figure.\n"
                "6. For every figure, change the amount of Go Fronts to fit the size requested.\n"
                "7. Combine all commands in a single response separated by a comma.\n"
                "8. Distances and rotations must be rounded to the nearest multiple of 10 cm or 30 degrees.\n"
                "9. Do not put any other characters than the commands and the commas. Do not end the list of commands with a period.\n"
                "\n\n"
                "Square Instructions:\n"
                "1. For a square with side length n cm, move forward n/10 times (rounded to the nearest multiple of 10 cm) at each side.\n"
                "2. After each movement, turn right 3 times (to make a 90-degree turn).\n"
                "3. Repeat for all 4 sides of the square.\n"
                "Example: 'Go Front (repeat n/10 times), Turn Right, Turn Right, Turn Right, Go Front (repeat n/10 times), Turn Right, Turn Right, Turn Right, Go Front (repeat n/10 times), Turn Right, Turn Right, Turn Right, Go Front (repeat n/10 times)'\n"
                "\n\n"
                "Circle Instructions:\n"
                "1. For a circle with radius n cm, calculate the circumference using the formula C = 2 * pi * n, where n is the radius.\n"
                "2. The drone can only move forward in increments of 10 cm, so round the forward movement to the nearest 10 cm to fit the circumference.\n"
                "3. Divide the circle into 12 equal segments of 30 degrees each (to complete 360 degrees). For each segment, move forward by a distance proportional to the radius.\n"
                "4. The distance moved forward per segment can be calculated by dividing the circle's circumference by 12 (the number of 30-degree segments). Since the drone moves in 10 cm increments, round the distance to the nearest 10 cm.\n"
                "   - For example, if the circumference is 628 cm, the drone moves approximately 52 cm per segment (rounded to the nearest 10 cm, which is 50 cm). This would mean the drone moves forward 5 times (10 cm each) in each 30-degree turn.\n"
                "5. Repeat the movement and rotation steps 12 times to approximate the circle.\n"
                "6. Example: For a circle with radius 100 cm (circumference = 628 cm), the drone would move forward 5 times (10 cm per step) for each of the 12 segments, turning right 30 degrees after each movement. The sequence of commands would look like this: 'Go Front, Go Front, Go Front, Go Front, Go Front, Turn Right, Go Front, Go Front, Go Front, Go Front, Go Front, Turn Right, ...' (repeat for 12 steps).\n"
            ),
        },
        {"role": "user", "content": chatbox_text},
    ],
)
\end{lstlisting}

% \input{appendix2}
%%%%%%%%%%%%%%%%%%%%%%%%%%%%%%%%%%%

\singlespacing

\setlength{\bibitemsep}{1.5em}

%\ihead{}
%\printbibliography[title=\Literaturverzeichnis] 
%\printbibliography 
%\cleardoublepage

%\initializeBibliography
%%%%%%%%%%%%%%%%%%%%%%%%%%%%%%%%%%%

%%%%%%%%%%%%%%%%%%%%%%%%%%%%%%%%%%%
% INDEX
% @stud: ggf. Index auskommentieren, wenn nicht benötigt
%
% \addcontentsline{toc}{chapter}{Index}
% \printindex

\end{document}